\documentclass[11pt,a4paper]{article}

\usepackage[margin=2.5cm]{geometry}
\usepackage[T1]{fontenc}
\usepackage{lmodern}
\usepackage{microtype}
\usepackage{graphicx}
\usepackage{xcolor}
\usepackage[most]{tcolorbox}
\usepackage{titlesec}
\usepackage{fancyhdr}
\usepackage{enumitem}
\usepackage{booktabs}
\usepackage{amsmath}
\usepackage{listings}
\usepackage{parskip}
\usepackage{hyperref}

% ── Colors ────────────────────────────────────────────────
\definecolor{black}{HTML}{000000}
\definecolor{darkgray}{HTML}{333333}
\definecolor{midgray}{HTML}{666666}
\definecolor{lightgray}{HTML}{F0F0F0}
\definecolor{rulegray}{HTML}{CCCCCC}
\definecolor{kwblue}{HTML}{003399}
\definecolor{strgreen}{HTML}{006633}
\definecolor{cmtgray}{HTML}{999999}

\hypersetup{
    colorlinks=true,
    linkcolor=darkgray,
    urlcolor=darkgray,
    citecolor=darkgray,
    pdftitle={Operating Systems Lab -- Assignment 2},
    pdfauthor={Ritam Ghosh},
}

% ── Header / Footer ──────────────────────────────────────
\pagestyle{fancy}
\fancyhf{}
\fancyhead[L]{\small\textsc{Operating Systems Lab --- Assignment 2}}
\fancyhead[R]{\small\textsc{Ritam Ghosh \textbar\ 002410501040}}
\fancyfoot[C]{\small\thepage}
\renewcommand{\headrulewidth}{0.4pt}
\renewcommand{\footrulewidth}{0.4pt}

% ── Section Formatting ────────────────────────────────────
\titleformat{\section}
  {\large\bfseries\scshape}{\thesection.}{0.6em}{}
  [\vspace{0.3em}{\color{rulegray}\rule{\linewidth}{0.6pt}}\vspace{0.5em}]

\titleformat{\subsection}
  {\normalsize\bfseries}{\thesubsection}{0.6em}{}

% ── Code Listings ─────────────────────────────────────────
\lstdefinestyle{cstyle}{
    language=C,
    basicstyle=\ttfamily\scriptsize,
    backgroundcolor=\color{lightgray},
    frame=single,
    rulecolor=\color{rulegray},
    keywordstyle=\color{kwblue}\bfseries,
    stringstyle=\color{strgreen},
    commentstyle=\color{cmtgray}\itshape,
    numbers=left,
    numberstyle=\tiny\color{cmtgray},
    numbersep=8pt,
    breaklines=true,
    breakatwhitespace=true,
    tabsize=4,
    showstringspaces=false,
    columns=flexible,
    keepspaces=true,
    xleftmargin=1.2em,
    xrightmargin=0.5em,
    aboveskip=0.6em,
    belowskip=0.6em,
    morekeywords={pid_t, fork, pipe, read, write, close, usleep, wait,
                  waitpid, printf, perror, exit, srand, rand, getpid,
                  fflush, mkfifo, open, unlink, msgsnd, msgrcv, msgget,
                  msgctl, ftok, strncpy, snprintf, atoi, strtok, fgets,
                  IPC_CREAT, IPC_RMID, O_RDONLY, O_WRONLY, NULL, STDIN_FILENO}
}

\lstdefinestyle{cppstyle}{
    language=C++,
    basicstyle=\ttfamily\scriptsize,
    backgroundcolor=\color{lightgray},
    frame=single,
    rulecolor=\color{rulegray},
    keywordstyle=\color{kwblue}\bfseries,
    stringstyle=\color{strgreen},
    commentstyle=\color{cmtgray}\itshape,
    numbers=left,
    numberstyle=\tiny\color{cmtgray},
    numbersep=8pt,
    breaklines=true,
    breakatwhitespace=true,
    tabsize=4,
    showstringspaces=false,
    columns=flexible,
    keepspaces=true,
    xleftmargin=1.2em,
    xrightmargin=0.5em,
    aboveskip=0.6em,
    belowskip=0.6em,
    morekeywords={vector, queue, priority_queue, string, pair, ifstream,
                  ofstream, cout, cerr, endl, string, to_string, sort,
                  min, max, fixed, setprecision, setw, left, right, INT_MAX}
}

\lstdefinestyle{outputstyle}{
    basicstyle=\ttfamily\scriptsize,
    backgroundcolor=\color{lightgray},
    frame=single,
    rulecolor=\color{rulegray},
    numbers=none,
    breaklines=true,
    breakatwhitespace=true,
    tabsize=4,
    showstringspaces=false,
    columns=flexible,
    keepspaces=true,
    xleftmargin=0.8em,
    xrightmargin=0.5em,
    aboveskip=0.4em,
    belowskip=0.4em,
}

% ── Boxes ─────────────────────────────────────────────────
\newtcolorbox{questionbox}{
    colback=white,
    colframe=darkgray,
    fonttitle=\bfseries,
    title={Question},
    arc=0mm,
    boxrule=0.5pt,
    left=4pt, right=4pt, top=2pt, bottom=2pt,
    breakable,
}

\newtcolorbox{theorybox}{
    colback=white,
    colframe=darkgray,
    coltitle=white,
    fonttitle=\bfseries,
    title={Theory},
    arc=0mm,
    boxrule=0.5pt,
    left=4pt, right=4pt, top=2pt, bottom=2pt,
    breakable,
}

\newtcolorbox{outputbox}[1][]{
    colback=white,
    colframe=rulegray,
    fonttitle=\bfseries\small,
    title={Output},
    arc=0mm,
    boxrule=0.4pt,
    left=4pt, right=4pt, top=2pt, bottom=2pt,
    breakable,
    #1
}

\newtcolorbox{codebox}[1][]{
    colback=white,
    colframe=rulegray,
    arc=0mm,
    boxrule=0.4pt,
    left=2pt, right=2pt, top=2pt, bottom=2pt,
    breakable,
    #1
}

% ── Title Page ────────────────────────────────────────────
\title{}
\date{}
\author{}

\begin{document}

% ── Custom Title Page ─────────────────────────────────────
\thispagestyle{empty}
\begin{center}
\vspace*{1cm}

{\large\textsc{Jadavpur University}}\\[0.2cm]
{\normalsize\textsc{Kolkata, West Bengal}}\\[0.8cm]

{\color{rulegray}\rule{0.7\linewidth}{1pt}}\\[0.8cm]

{\LARGE\bfseries\textsc{Operating Systems Laboratory}}\\[0.4cm]
{\Large\textsc{Assignment --- II}}\\[0.8cm]

{\color{rulegray}\rule{0.7\linewidth}{1pt}}\\[1.2cm]

\begin{tabular}{rl}
\textbf{Name} & Ritam Ghosh \\[0.3cm]
\textbf{Roll No.} & 002410501040 \\[0.3cm]
\textbf{Section} & BCSE3.1 \\[0.3cm]
\textbf{Semester} & 3 (2026--27) \\[0.3cm]
\textbf{Subject} & Operating Systems Lab (CS39002) \\
\end{tabular}

\vfill

{\color{rulegray}\rule{0.7\linewidth}{1pt}}\\[0.3cm]
{\small\today}
\end{center}

\newpage
\setcounter{page}{1}

% ══════════════════════════════════════════════════════════
%  QUESTION 1
% ══════════════════════════════════════════════════════════
\section{IPC Synchronization using Pipes}

\begin{questionbox}
Two processes, Process X and Process Y, interact with data $x$ and $y$. Process X writes $x$ and reads $y$; Process Y writes $y$ and reads $x$. Ensure proper synchronization using \texttt{fork()} and pipes.
\end{questionbox}

\begin{theorybox}
A \textbf{pipe} is a unidirectional IPC channel with a write end and a read end. Its key property is \emph{blocking}: a \texttt{read()} on an empty pipe suspends the calling process until data is written.

Two pipes are used for bidirectional flow:
\begin{itemize}[nosep]
    \item Pipe 1 ($x$ flow): X writes, Y reads.
    \item Pipe 2 ($y$ flow): Y writes, X reads.
\end{itemize}

Both processes follow a \texttt{Write} then \texttt{Read} pattern. If the peer has not yet written, the read blocks, providing synchronization without explicit locks.
\end{theorybox}

\subsection*{Source Code}
\begin{codebox}
\begin{lstlisting}[style=cstyle]
#include <stdio.h>
#include <stdlib.h>
#include <sys/types.h>
#include <sys/wait.h>
#include <unistd.h>

int main() {
  int pipe1[2], pipe2[2];
  if (pipe(pipe1) == -1 || pipe(pipe2) == -1) {
    perror("Pipe creation failed");
    return 1;
  }

  pid_t pid = fork();
  if (pid > 0) {
    // Process X (Parent)
    close(pipe1[0]);
    close(pipe2[1]);

    int x = 100;
    printf("Process X is writing %d\n", x);
    write(pipe1[1], &x, sizeof(x));

    int y_read;
    printf("Process X is waiting to read \n");
    read(pipe2[0], &y_read, sizeof(y_read));
    printf("Process X read %d\n", y_read);

    close(pipe1[1]);
    close(pipe2[0]);
  } else {
    // Process Y (Child)
    close(pipe1[1]);
    close(pipe2[0]);

    int y = 200;
    printf("Process Y is writing %d\n", y);
    write(pipe2[1], &y, sizeof(y));

    int x_read;
    printf("Process Y is waiting to read \n");
    read(pipe1[0], &x_read, sizeof(x_read));
    printf("Process Y read %d\n", x_read);

    close(pipe1[0]);
    close(pipe2[1]);
  }
  return 0;
}
\end{lstlisting}
\end{codebox}

\subsection*{Output}
\begin{outputbox}
\begin{lstlisting}[style=outputstyle]
Process X is writing 100
Process X is waiting to read
Process X read 200
Process Y is writing 200
Process Y is waiting to read
Process Y read 100
\end{lstlisting}
\end{outputbox}

\newpage

% ══════════════════════════════════════════════════════════
%  QUESTION 2
% ══════════════════════════════════════════════════════════
\section{Process Interleaving and Strict Synchronization}

\begin{questionbox}
\textbf{(a)} Create two child processes X and Y that each run $N$ iterations, printing their identity and iteration number with random sleep.\\
\textbf{(b)} Enforce strict alternation: Y cannot start iteration $i$ until X finishes iteration $i-1$.
\end{questionbox}

\begin{theorybox}
\textbf{Process Interleaving.} The OS scheduler uses time-slicing to multiplex processes. Random \texttt{usleep()} calls force context switches, producing non-deterministic output order each run.

\textbf{Strict Synchronization.} A pipe carries a token from X to Y. After each iteration, X writes a token; at the start of each iteration, Y reads it. The blocking property of \texttt{read()} ensures Y waits until X finishes.
\end{theorybox}

\subsection*{Part (a) --- \texttt{qs2a.c}}
\begin{codebox}
\begin{lstlisting}[style=cstyle]
#include <stdio.h>
#include <stdlib.h>
#include <unistd.h>
#include <sys/types.h>
#include <sys/wait.h>
#include <time.h>

#define ITERATIONS 10

void run_child(const char* name) {
    srand(time(NULL) ^ getpid());
    for (int i = 1; i <= ITERATIONS; i++) {
        printf("Process %s - Iteration %d\n", name, i);
        fflush(stdout);
        usleep(rand() % 100000);
    }
    exit(0);
}

int main() {
    pid_t pid_x = fork();
    if (pid_x < 0) { perror("Fork failed for X"); exit(1); }
    else if (pid_x == 0) { run_child("X"); }

    pid_t pid_y = fork();
    if (pid_y < 0) { perror("Fork failed for Y"); exit(1); }
    else if (pid_y == 0) { run_child("Y"); }

    wait(NULL);
    wait(NULL);
    printf("Both processes finished.\n");
    return 0;
}
\end{lstlisting}
\end{codebox}

\subsection*{Output --- Part (a)}
\begin{outputbox}
\begin{lstlisting}[style=outputstyle]
Process X - Iteration 1
Process Y - Iteration 1
Process Y - Iteration 2
Process Y - Iteration 3
Process X - Iteration 2
Process Y - Iteration 4
Process Y - Iteration 5
Process X - Iteration 3
Process X - Iteration 4
Process Y - Iteration 6
Process X - Iteration 5
Process Y - Iteration 7
Process X - Iteration 6
Process X - Iteration 7
Process Y - Iteration 8
Process X - Iteration 8
Process Y - Iteration 9
Process Y - Iteration 10
Process X - Iteration 9
Process X - Iteration 10
Both processes finished.
\end{lstlisting}
\end{outputbox}

\subsection*{Part (b) --- \texttt{qs2b.c}}
\begin{codebox}
\begin{lstlisting}[style=cstyle]
#include <stdio.h>
#include <stdlib.h>
#include <sys/types.h>
#include <sys/wait.h>
#include <time.h>
#include <unistd.h>

#define ITERATIONS 10

int main() {
  int pipe_xy[2], pipe_yx[2];
  if (pipe(pipe_xy) == -1 || pipe(pipe_yx) == -1) {
    perror("Pipe creation failed");
    return 1;
  }

  pid_t pid = fork();
  if (pid < 0) { perror("Fork failed"); exit(1); }

  srand(time(NULL) ^ getpid());
  int token = 1;

  if (pid > 0) {
    // PROCESS X (Parent)
    close(pipe_xy[0]);
    close(pipe_yx[1]);

    for (int i = 1; i <= ITERATIONS; i++) {
      printf("Process X - Iteration %d\n", i);
      fflush(stdout);
      usleep(rand() % 100000);
      write(pipe_xy[1], &token, sizeof(token));
      if (i < ITERATIONS) {
        read(pipe_yx[0], &token, sizeof(token));
      }
    }
    close(pipe_xy[1]);
    close(pipe_yx[0]);
    wait(NULL);
    printf("Both processes finished.\n");
  } else {
    // PROCESS Y (Child)
    close(pipe_xy[1]);
    close(pipe_yx[0]);

    for (int i = 1; i <= ITERATIONS; i++) {
      read(pipe_xy[0], &token, sizeof(token));
      printf("Process Y - Iteration %d\n", i);
      fflush(stdout);
      usleep(rand() % 100000);
      if (i < ITERATIONS) {
        write(pipe_yx[1], &token, sizeof(token));
      }
    }
    close(pipe_xy[0]);
    close(pipe_yx[1]);
    exit(0);
  }
  return 0;
}
\end{lstlisting}
\end{codebox}

\subsection*{Output --- Part (b)}
\begin{outputbox}
\begin{lstlisting}[style=outputstyle]
Process X - Iteration 1
Process Y - Iteration 1
Process X - Iteration 2
Process Y - Iteration 2
Process X - Iteration 3
Process Y - Iteration 3
Process X - Iteration 4
Process Y - Iteration 4
Process X - Iteration 5
Process Y - Iteration 5
Process X - Iteration 6
Process Y - Iteration 6
Process X - Iteration 7
Process Y - Iteration 7
Process X - Iteration 8
Process Y - Iteration 8
Process X - Iteration 9
Process Y - Iteration 9
Process X - Iteration 10
Process Y - Iteration 10
Both processes finished.
\end{lstlisting}
\end{outputbox}

\newpage

% ══════════════════════════════════════════════════════════
%  QUESTION 3
% ══════════════════════════════════════════════════════════
\section{IPC Mechanisms: Pipes, FIFOs, and Message Queues}

\begin{questionbox}
Implement three IPC applications:
\begin{enumerate}[nosep]
    \item[(a)] Weather broadcast using \textbf{Message Queues} (System V).
    \item[(b)] Telephonic conversation using \textbf{unnamed Pipes}.
    \item[(c)] Broadcast with listener replies using \textbf{named FIFOs}.
\end{enumerate}
\end{questionbox}

\begin{theorybox}
\begin{itemize}[nosep]
    \item \textbf{Pipes:} Unidirectional byte streams between related processes. Two pipes are needed for bidirectional communication.
    \item \textbf{FIFOs (Named Pipes):} Filesystem entries that allow IPC between unrelated processes via \texttt{mkfifo()}.
    \item \textbf{Message Queues (System V):} Kernel-managed linked lists of typed messages. The \texttt{mtype} field enables selective reception by different listeners.
\end{itemize}
\end{theorybox}

\subsection*{Part (a) --- Weather Broadcast (\texttt{qs3a.c})}
\begin{codebox}
\begin{lstlisting}[style=cstyle]
#include <stdio.h>
#include <stdlib.h>
#include <string.h>
#include <sys/ipc.h>
#include <sys/msg.h>
#include <sys/types.h>
#include <sys/wait.h>
#include <unistd.h>

#define NUM_LISTENERS 3
#define NUM_BROADCASTS 5
#define MSG_SIZE 128

struct msgbuf {
  long mtype;
  char mtext[MSG_SIZE];
};

void listener(int msgid, int listener_id) {
  struct msgbuf msg;
  long my_type = listener_id + 1;
  printf("  [Listener %d] Started (PID: %d), listening on mtype=%ld\n",
         listener_id, getpid(), my_type);
  fflush(stdout);
  while (1) {
    if (msgrcv(msgid, &msg, MSG_SIZE, my_type, 0) == -1) {
      perror("  msgrcv"); exit(1);
    }
    if (strcmp(msg.mtext, "END") == 0) {
      printf("  [Listener %d] Received END. Exiting.\n", listener_id);
      fflush(stdout); break;
    }
    printf("  [Listener %d] Received: \"%s\"\n", listener_id, msg.mtext);
    fflush(stdout);
  }
  exit(0);
}

int main() {
  printf("========================================\n");
  printf("  3a: Weather Broadcast (Message Queue)\n");
  printf("========================================\n\n");
  key_t key = ftok("/tmp", 'W');
  if (key == -1) { perror("ftok"); return 1; }
  int msgid = msgget(key, IPC_CREAT | 0666);
  if (msgid == -1) { perror("msgget"); return 1; }
  printf("  [Broadcaster] Message queue created (id=%d)\n\n", msgid);
  fflush(stdout);

  pid_t pids[NUM_LISTENERS];
  for (int i = 0; i < NUM_LISTENERS; i++) {
    pids[i] = fork();
    if (pids[i] < 0) { perror("fork"); return 1; }
    if (pids[i] == 0) { listener(msgid, i); exit(0); }
  }
  usleep(100000);

  const char *weather[] = {"Sunny, 32C, Humidity 45%",
    "Partly Cloudy, 28C, Wind 15km/h", "Thunderstorm Warning, 24C",
    "Clear Skies, 30C, UV Index High", "Light Rain, 22C, Humidity 80%"};

  struct msgbuf msg;
  for (int b = 0; b < NUM_BROADCASTS; b++) {
    printf("\n  [Broadcaster] Sending: \"%s\"\n", weather[b]);
    fflush(stdout);
    for (int i = 0; i < NUM_LISTENERS; i++) {
      msg.mtype = i + 1;
      strncpy(msg.mtext, weather[b], MSG_SIZE - 1);
      msg.mtext[MSG_SIZE - 1] = '\0';
      if (msgsnd(msgid, &msg, strlen(msg.mtext) + 1, 0) == -1) {
        perror("  msgsnd"); return 1;
      }
    }
    usleep(200000);
  }

  printf("\n  [Broadcaster] Sending END to all.\n");
  fflush(stdout);
  for (int i = 0; i < NUM_LISTENERS; i++) {
    msg.mtype = i + 1;
    strncpy(msg.mtext, "END", MSG_SIZE);
    if (msgsnd(msgid, &msg, strlen(msg.mtext) + 1, 0) == -1)
      perror("  msgsnd END");
  }
  for (int i = 0; i < NUM_LISTENERS; i++) waitpid(pids[i], NULL, 0);
  if (msgctl(msgid, IPC_RMID, NULL) == -1) perror("msgctl IPC_RMID");
  else printf("\n  [Broadcaster] Queue removed. Done.\n");
  return 0;
}
\end{lstlisting}
\end{codebox}

\subsection*{Output --- Part (a)}
\begin{outputbox}
\begin{lstlisting}[style=outputstyle]
========================================
  3a: Weather Broadcast (Message Queue)
========================================

  [Broadcaster] Message queue created (id=65536)

  [Listener 0] Started (PID: 52047), listening on mtype=1
  [Listener 1] Started (PID: 52048), listening on mtype=2
  [Listener 2] Started (PID: 52049), listening on mtype=3

  [Broadcaster] Sending: "Sunny, 32C, Humidity 45%"
  [Listener 2] Received: "Sunny, 32C, Humidity 45%"
  [Listener 0] Received: "Sunny, 32C, Humidity 45%"
  [Listener 1] Received: "Sunny, 32C, Humidity 45%"

  [Broadcaster] Sending: "Partly Cloudy, 28C, Wind 15km/h"
  [Listener 2] Received: "Partly Cloudy, 28C, Wind 15km/h"
  [Listener 0] Received: "Partly Cloudy, 28C, Wind 15km/h"
  [Listener 1] Received: "Partly Cloudy, 28C, Wind 15km/h"

  [Broadcaster] Sending: "Thunderstorm Warning, 24C"
  [Listener 2] Received: "Thunderstorm Warning, 24C"
  [Listener 0] Received: "Thunderstorm Warning, 24C"
  [Listener 1] Received: "Thunderstorm Warning, 24C"

  [Broadcaster] Sending: "Clear Skies, 30C, UV Index High"
  [Listener 1] Received: "Clear Skies, 30C, UV Index High"
  [Listener 0] Received: "Clear Skies, 30C, UV Index High"
  [Listener 2] Received: "Clear Skies, 30C, UV Index High"

  [Broadcaster] Sending: "Light Rain, 22C, Humidity 80%"
  [Listener 0] Received: "Light Rain, 22C, Humidity 80%"
  [Listener 2] Received: "Light Rain, 22C, Humidity 80%"
  [Listener 1] Received: "Light Rain, 22C, Humidity 80%"

  [Broadcaster] Sending END signal to all listeners.
  [Listener 2] Received END signal. Exiting.
  [Listener 0] Received END signal. Exiting.
  [Listener 1] Received END signal. Exiting.

  [Broadcaster] Message queue removed. Done.
\end{lstlisting}
\end{outputbox}

\newpage

\subsection*{Part (b) --- Telephonic Conversation (\texttt{qs3b.c})}
\begin{codebox}
\begin{lstlisting}[style=cstyle]
#include <stdio.h>
#include <stdlib.h>
#include <string.h>
#include <sys/types.h>
#include <sys/wait.h>
#include <unistd.h>

#define MSG_SIZE 256
#define NUM_EXCHANGES 4

int main() {
  printf("========================================\n");
  printf("  3b: Telephonic Conversation (Pipe)\n");
  printf("========================================\n\n");
  fflush(stdout);

  int pipe_c2r[2], pipe_r2c[2];
  if (pipe(pipe_c2r) == -1 || pipe(pipe_r2c) == -1) {
    perror("pipe"); return 1;
  }

  const char *caller_lines[] = {"Hello? Is this the Domino's Pizza?",
    "I'd like to order a large pepperoni pizza.",
    "Yes, with extra cheese please.", "Great, thanks! Goodbye."};
  const char *receiver_lines[] = {
    "Yes, this is Domino's! How can I help you?",
    "Sure! Any toppings you'd like to add?",
    "Got it! Your order will be ready in 30 minutes.",
    "Thank you for calling! Goodbye."};

  pid_t pid = fork();
  if (pid < 0) { perror("fork"); return 1; }

  char buf[MSG_SIZE];
  if (pid > 0) {
    // CALLER (Parent)
    close(pipe_c2r[0]); close(pipe_r2c[1]);
    for (int i = 0; i < NUM_EXCHANGES; i++) {
      printf("  [Caller]   : %s\n", caller_lines[i]); fflush(stdout);
      write(pipe_c2r[1], caller_lines[i], strlen(caller_lines[i]) + 1);
      read(pipe_r2c[0], buf, MSG_SIZE);
      printf("  [Caller]   hears: \"%s\"\n", buf); fflush(stdout);
      usleep(300000);
    }
    close(pipe_c2r[1]); close(pipe_r2c[0]);
    wait(NULL);
    printf("\n  --- Call ended ---\n");
  } else {
    // RECEIVER (Child)
    close(pipe_c2r[1]); close(pipe_r2c[0]);
    for (int i = 0; i < NUM_EXCHANGES; i++) {
      read(pipe_c2r[0], buf, MSG_SIZE);
      printf("  [Receiver] hears: \"%s\"\n", buf); fflush(stdout);
      usleep(200000);
      printf("  [Receiver] : %s\n", receiver_lines[i]); fflush(stdout);
      write(pipe_r2c[1], receiver_lines[i], strlen(receiver_lines[i]) + 1);
    }
    close(pipe_c2r[0]); close(pipe_r2c[1]);
    exit(0);
  }
  return 0;
}
\end{lstlisting}
\end{codebox}

\subsection*{Output --- Part (b)}
\begin{outputbox}
\begin{lstlisting}[style=outputstyle]
========================================
  3b: Telephonic Conversation (Pipe)
========================================

  [Caller]   : Hello? Is this the Domino's Pizza?
  [Receiver] hears: "Hello? Is this the Domino's Pizza?"
  [Receiver] : Yes, this is Domino's! How can I help you?
  [Caller]   hears: "Yes, this is Domino's! How can I help you?"
  [Caller]   : I'd like to order a large pepperoni pizza.
  [Receiver] hears: "I'd like to order a large pepperoni pizza."
  [Receiver] : Sure! Any toppings you'd like to add?
  [Caller]   hears: "Sure! Any toppings you'd like to add?"
  [Caller]   : Yes, with extra cheese please.
  [Receiver] hears: "Yes, with extra cheese please."
  [Receiver] : Got it! Your order will be ready in 30 minutes.
  [Caller]   hears: "Got it! Your order will be ready in 30 minutes."
  [Caller]   : Great, thanks! Goodbye.
  [Receiver] hears: "Great, thanks! Goodbye."
  [Receiver] : Thank you for calling! Goodbye.
  [Caller]   hears: "Thank you for calling! Goodbye."

  --- Call ended ---
\end{lstlisting}
\end{outputbox}

\newpage

\subsection*{Part (c) --- Broadcast + Replies via FIFOs (\texttt{qs3c.c})}
\begin{codebox}
\begin{lstlisting}[style=cstyle]
#include <errno.h>
#include <fcntl.h>
#include <stdio.h>
#include <stdlib.h>
#include <string.h>
#include <sys/stat.h>
#include <sys/types.h>
#include <sys/wait.h>
#include <unistd.h>

#define NUM_LISTENERS 3
#define NUM_BROADCASTS 4
#define MSG_SIZE 256

void get_bcast_path(int id, char *buf) {
  sprintf(buf, "/tmp/qs3_bcast_%d", id);
}
void get_reply_path(int id, char *buf) {
  sprintf(buf, "/tmp/qs3_reply_%d", id);
}

void listener(int id) {
  char bcast_path[64], reply_path[64], buf[MSG_SIZE];
  get_bcast_path(id, bcast_path);
  get_reply_path(id, reply_path);
  printf("  [Listener %d] Started (PID: %d)\n", id, getpid());
  fflush(stdout);

  int bcast_fd = open(bcast_path, O_RDONLY);
  if (bcast_fd == -1) { perror("  listener open bcast"); exit(1); }
  int reply_fd = open(reply_path, O_WRONLY);
  if (reply_fd == -1) { perror("  listener open reply"); exit(1); }

  while (1) {
    ssize_t n = read(bcast_fd, buf, MSG_SIZE);
    if (n <= 0) break;
    if (strcmp(buf, "END") == 0) {
      printf("  [Listener %d] Received END. Exiting.\n", id);
      fflush(stdout);
      char ack[MSG_SIZE];
      snprintf(ack, MSG_SIZE, "Listener %d signing off.", id);
      write(reply_fd, ack, strlen(ack) + 1);
      break;
    }
    printf("  [Listener %d] Received: \"%s\"\n", id, buf);
    fflush(stdout);
    char reply[MSG_SIZE];
    snprintf(reply, MSG_SIZE, "Listener %d: Acknowledged \"%s\"", id, buf);
    write(reply_fd, reply, strlen(reply) + 1);
  }
  close(bcast_fd); close(reply_fd); exit(0);
}

int main() {
  printf("========================================\n");
  printf("  3c: Broadcast + Replies (Named FIFO)\n");
  printf("========================================\n\n"); fflush(stdout);

  char bcast_paths[NUM_LISTENERS][64], reply_paths[NUM_LISTENERS][64];
  for (int i = 0; i < NUM_LISTENERS; i++) {
    get_bcast_path(i, bcast_paths[i]);
    get_reply_path(i, reply_paths[i]);
    unlink(bcast_paths[i]); unlink(reply_paths[i]);
    if (mkfifo(bcast_paths[i], 0666) == -1) { perror("mkfifo"); return 1; }
    if (mkfifo(reply_paths[i], 0666) == -1) { perror("mkfifo"); return 1; }
  }
  printf("  [Broadcaster] FIFOs created.\n\n"); fflush(stdout);

  pid_t pids[NUM_LISTENERS];
  for (int i = 0; i < NUM_LISTENERS; i++) {
    pids[i] = fork();
    if (pids[i] < 0) { perror("fork"); return 1; }
    if (pids[i] == 0) { listener(i); exit(0); }
  }

  int bcast_fds[NUM_LISTENERS], reply_fds[NUM_LISTENERS];
  for (int i = 0; i < NUM_LISTENERS; i++) {
    bcast_fds[i] = open(bcast_paths[i], O_WRONLY);
    if (bcast_fds[i] == -1) { perror("open bcast"); return 1; }
  }
  for (int i = 0; i < NUM_LISTENERS; i++) {
    reply_fds[i] = open(reply_paths[i], O_RDONLY);
    if (reply_fds[i] == -1) { perror("open reply"); return 1; }
  }

  const char *updates[] = {"Heatwave alert: 42C expected tomorrow",
    "Rain forecast for evening, carry umbrellas",
    "Air quality index: Moderate (AQI 120)",
    "Weekend outlook: Pleasant, 26C"};
  char buf[MSG_SIZE];

  for (int b = 0; b < NUM_BROADCASTS; b++) {
    printf("\n  [Broadcaster] Sending: \"%s\"\n", updates[b]);
    fflush(stdout);
    for (int i = 0; i < NUM_LISTENERS; i++)
      write(bcast_fds[i], updates[b], strlen(updates[b]) + 1);
    usleep(100000);
    for (int i = 0; i < NUM_LISTENERS; i++) {
      ssize_t n = read(reply_fds[i], buf, MSG_SIZE);
      if (n > 0) {
        printf("  [Broadcaster] Reply: \"%s\"\n", buf);
        fflush(stdout);
      }
    }
    usleep(200000);
  }

  printf("\n  [Broadcaster] Sending END.\n"); fflush(stdout);
  for (int i = 0; i < NUM_LISTENERS; i++) write(bcast_fds[i], "END", 4);
  usleep(100000);
  for (int i = 0; i < NUM_LISTENERS; i++) {
    ssize_t n = read(reply_fds[i], buf, MSG_SIZE);
    if (n > 0) {
      printf("  [Broadcaster] Final: \"%s\"\n", buf);
      fflush(stdout);
    }
  }
  for (int i = 0; i < NUM_LISTENERS; i++)
    { close(bcast_fds[i]); close(reply_fds[i]); }
  for (int i = 0; i < NUM_LISTENERS; i++) waitpid(pids[i], NULL, 0);
  for (int i = 0; i < NUM_LISTENERS; i++)
    { unlink(bcast_paths[i]); unlink(reply_paths[i]); }
  printf("\n  [Broadcaster] FIFOs cleaned up. Done.\n");
  return 0;
}
\end{lstlisting}
\end{codebox}

\subsection*{Output --- Part (c)}
\begin{outputbox}
\begin{lstlisting}[style=outputstyle]
========================================
  3c: Broadcast + Replies (Named FIFO)
========================================

  [Broadcaster] FIFOs created.

  [Listener 0] Started (PID: 52053)
  [Listener 1] Started (PID: 52054)
  [Listener 2] Started (PID: 52055)

  [Broadcaster] Sending: "Heatwave alert: 42C expected tomorrow"
  [Listener 0] Received: "Heatwave alert: 42C expected tomorrow"
  [Listener 1] Received: "Heatwave alert: 42C expected tomorrow"
  [Listener 2] Received: "Heatwave alert: 42C expected tomorrow"
  [Broadcaster] Reply: "Listener 0: Acknowledged ..."
  [Broadcaster] Reply: "Listener 1: Acknowledged ..."
  [Broadcaster] Reply: "Listener 2: Acknowledged ..."

  [Broadcaster] Sending: "Rain forecast for evening, carry umbrellas"
  [Listener 0] Received: "Rain forecast for evening, carry umbrellas"
  [Listener 1] Received: "Rain forecast for evening, carry umbrellas"
  [Listener 2] Received: "Rain forecast for evening, carry umbrellas"
  [Broadcaster] Reply: "Listener 0: Acknowledged ..."
  [Broadcaster] Reply: "Listener 1: Acknowledged ..."
  [Broadcaster] Reply: "Listener 2: Acknowledged ..."

  [Broadcaster] Sending: "Air quality index: Moderate (AQI 120)"
  [Listener 0] Received: "Air quality index: Moderate (AQI 120)"
  [Listener 1] Received: "Air quality index: Moderate (AQI 120)"
  [Listener 2] Received: "Air quality index: Moderate (AQI 120)"
  [Broadcaster] Reply: "Listener 0: Acknowledged ..."
  [Broadcaster] Reply: "Listener 1: Acknowledged ..."
  [Broadcaster] Reply: "Listener 2: Acknowledged ..."

  [Broadcaster] Sending: "Weekend outlook: Pleasant, 26C"
  [Listener 0] Received: "Weekend outlook: Pleasant, 26C"
  [Listener 2] Received: "Weekend outlook: Pleasant, 26C"
  [Listener 1] Received: "Weekend outlook: Pleasant, 26C"
  [Broadcaster] Reply: "Listener 0: Acknowledged ..."
  [Broadcaster] Reply: "Listener 1: Acknowledged ..."
  [Broadcaster] Reply: "Listener 2: Acknowledged ..."

  [Broadcaster] Sending END to all listeners.
  [Listener 0] Received END. Exiting.
  [Listener 1] Received END. Exiting.
  [Listener 2] Received END. Exiting.
  [Broadcaster] Final: "Listener 0 signing off."
  [Broadcaster] Final: "Listener 1 signing off."
  [Broadcaster] Final: "Listener 2 signing off."

  [Broadcaster] FIFOs cleaned up. Done.
\end{lstlisting}
\end{outputbox}

\newpage

% ══════════════════════════════════════════════════════════
%  QUESTION 4
% ══════════════════════════════════════════════════════════
\section{CPU Scheduling Algorithms}

\begin{questionbox}
Build a scheduler that reads job execution profiles from a file and simulates \textbf{FCFS}, \textbf{Non-Preemptive Priority}, and \textbf{Round Robin (quantum = 16)}. Compare the average turnaround time (TAT) and waiting time (WT).
\end{questionbox}

\begin{theorybox}
A \textbf{CPU Scheduler} selects which ready process gets the CPU next.
\begin{itemize}[nosep]
    \item \textbf{FCFS} --- Non-preemptive, FIFO order. Simple but suffers from the convoy effect.
    \item \textbf{Priority} --- Non-preemptive, lowest priority number wins. Risk of starvation for low-priority jobs.
    \item \textbf{Round Robin} --- Preemptive, fixed time quantum. Fair but incurs context-switch overhead.
\end{itemize}
Each job has alternating CPU and I/O bursts. After a CPU burst, the job enters I/O; after I/O completes, it re-enters the ready queue. Waiting time accumulates only while in the ready queue.
\end{theorybox}

\subsection*{Input File --- \texttt{jobs.txt}}
\begin{codebox}
\begin{lstlisting}[style=outputstyle]
1 3 0 32 64 16 32 -1
2 1 4 64 8 32 16 8 -1
3 4 6 24 48 12 -1
4 2 8 48 32 24 64 16 -1
5 5 10 16 32 48 -1
6 1 12 64 16 32 8 16 -1
7 3 14 32 24 16 48 8 -1
8 2 16 24 16 48 32 16 -1
\end{lstlisting}
\end{codebox}

\subsection*{Output}
\begin{outputbox}
\begin{lstlisting}[style=outputstyle]
========================================
     CPU Scheduler Simulation
========================================

  Loaded 8 job profiles:
  ------------------------------------------------------------
  Job 1 | Priority: 3 | Arrival:  0 | CPU(32)->I/O(64)->CPU(16)->I/O(32)
  Job 2 | Priority: 1 | Arrival:  4 | CPU(64)->I/O(8)->CPU(32)->I/O(16)->CPU(8)
  Job 3 | Priority: 4 | Arrival:  6 | CPU(24)->I/O(48)->CPU(12)
  Job 4 | Priority: 2 | Arrival:  8 | CPU(48)->I/O(32)->CPU(24)->I/O(64)->CPU(16)
  Job 5 | Priority: 5 | Arrival: 10 | CPU(16)->I/O(32)->CPU(48)
  Job 6 | Priority: 1 | Arrival: 12 | CPU(64)->I/O(16)->CPU(32)->I/O(8)->CPU(16)
  Job 7 | Priority: 3 | Arrival: 14 | CPU(32)->I/O(24)->CPU(16)->I/O(48)->CPU(8)
  Job 8 | Priority: 2 | Arrival: 16 | CPU(24)->I/O(16)->CPU(48)->I/O(32)->CPU(16)
  ------------------------------------------------------------

  ===== 1. FCFS Scheduling =====

  Gantt Chart:
  | J1  | J2  | J3  | J4  | J5  | J6  | J7  | J8  | J1  | ...
  0    32    96   120   168   184   248   280   304   320  ...

  +----------- FCFS (First Come First Serve) -------------+
  | Job   Arrival   Completion   TAT        WT             |
  | 1     0         352          352        208            |
  | 2     4         540          536        372            |
  | 3     6         364          358        274            |
  | 4     8         556          548        300            |
  | 5     10        436          426        298            |
  | 6     12        572          560        400            |
  | 7     14        580          566        438            |
  | 8     16        596          580        436            |
  | Average TAT:                490.75                     |
  | Average WT:                 340.75                     |
  +-------------------------------------------------------+

  ===== 2. Non-Preemptive Priority Scheduling =====

  +------------- Non-Preemptive Priority -----------------+
  | Job   Arrival   Completion   TAT        WT             |
  | 1     0         368          368        224            |
  | 2     4         232          228        28             |
  | 3     6         572          566        482            |
  | 4     8         456          448        240            |
  | 5     10        620          610        514            |
  | 6     12        248          236        84             |
  | 7     14        544          530        394            |
  | 8     16        472          456        320            |
  | Average TAT:                430.25                     |
  | Average WT:                 285.75                     |
  +-------------------------------------------------------+

  ===== 3. Round Robin (quantum=16) =====

  +--------------- Round Robin (q=16) --------------------+
  | Job   Arrival   Completion   TAT        WT             |
  | 1     0         336          336        184            |
  | 2     4         548          544        408            |
  | 3     6         316          310        226            |
  | 4     8         596          588        396            |
  | 5     10        428          418        322            |
  | 6     12        580          568        432            |
  | 7     14        468          454        326            |
  | 8     16        564          548        412            |
  | Average TAT:                470.75                     |
  | Average WT:                 338.25                     |
  +-------------------------------------------------------+
\end{lstlisting}
\end{outputbox}

\subsection*{Comparison}
\begin{table}[h]
\centering
\begin{tabular}{lcc}
\toprule
\textbf{Algorithm} & \textbf{Avg TAT} & \textbf{Avg WT} \\
\midrule
FCFS & 490.75 & 340.75 \\
Non-Preemptive Priority & \textbf{430.25} & \textbf{285.75} \\
Round Robin ($q = 16$) & 470.75 & 338.25 \\
\bottomrule
\end{tabular}
\end{table}

Priority scheduling achieves the best (lowest) average waiting time for this workload, while FCFS performs worst due to the convoy effect.

\newpage

% ══════════════════════════════════════════════════════════
%  QUESTION 5
% ══════════════════════════════════════════════════════════
\section{Deadlock Avoidance --- Banker's Algorithm}

\begin{questionbox}
Implement the Banker's Algorithm to determine if a system is in a safe state and whether resource requests should be granted without leading to deadlock. Simulate until all processes finish.
\end{questionbox}

\begin{theorybox}
The \textbf{Banker's Algorithm} avoids deadlock by checking safety before granting resources.

Key data structures:
\begin{itemize}[nosep]
    \item \textbf{Available} --- Free instances per resource type.
    \item \textbf{Max} --- Maximum demand of each process.
    \item \textbf{Allocation} --- Currently held resources.
    \item \textbf{Need} --- Remaining demand ($\text{Need} = \text{Max} - \text{Allocation}$).
\end{itemize}

\textbf{Safety check:} Find a process whose Need $\leq$ Work; simulate its completion (Work += Allocation); repeat. If all processes finish, the state is safe.

\textbf{Resource request:} Verify Request $\leq$ Need and Request $\leq$ Available; tentatively allocate; run safety check; grant if safe, rollback if unsafe.
\end{theorybox}

\subsection*{Input File --- \texttt{input.txt}}
\begin{codebox}
\begin{lstlisting}[style=outputstyle]
<#no of resources> 4 <#instances each> 2 4 5 3
<#no of processes> 3 <max needs> 1 1 1 1, 2 3 1 2, 2 2 1 3
\end{lstlisting}
\end{codebox}

\subsection*{Output}
\begin{outputbox}
\begin{lstlisting}[style=outputstyle]
========================================
   Banker's Algorithm Simulation
========================================
  Processes : 3
  Resources : 4
  Total     : [  2  4  5  3 ]

  +------- Current System State -------+
  Available : [  2  4  5  3 ]

           Allocation     Need
           R0 R1 R2 R3       R0 R1 R2 R3
    P0  [  0  0  0  0 ]   [  1  1  1  1 ]
    P1  [  0  0  0  0 ]   [  2  3  1  2 ]
    P2  [  0  0  0  0 ]   [  2  2  1  3 ]

  ---- Step 1 ----
  P2 requests : [  2  2  1  3 ]
  -> GRANT : Safe state confirmed.
     Safe sequence: < P2 P0 P1 >
  >> P2 completed. Releasing: [ 2 2 1 3 ]
     Execution order: < P2 >

  +------- Current System State -------+
  Available : [  2  4  5  3 ]

           Allocation     Need
           R0 R1 R2 R3       R0 R1 R2 R3
    P0  [  0  0  0  0 ]   [  1  1  1  1 ]
    P1  [  0  0  0  0 ]   [  2  3  1  2 ]
    P2* [  0  0  0  0 ]   [  0  0  0  0 ]

  ---- Step 2 ----
  P0 requests : [  1  1  1  1 ]
  -> GRANT : Safe state confirmed.
     Safe sequence: < P0 P1 P2 >
  >> P0 completed. Releasing: [ 1 1 1 1 ]
     Execution order: < P2 P0 >

  +------- Current System State -------+
  Available : [  2  4  5  3 ]

           Allocation     Need
           R0 R1 R2 R3       R0 R1 R2 R3
    P0* [  0  0  0  0 ]   [  0  0  0  0 ]
    P1  [  0  0  0  0 ]   [  2  3  1  2 ]
    P2* [  0  0  0  0 ]   [  0  0  0  0 ]

  ---- Step 3 ----
  P1 requests : [  2  3  1  2 ]
  -> GRANT : Safe state confirmed.
     Safe sequence: < P1 P2 P0 >
  >> P1 completed. Releasing: [ 2 3 1 2 ]
     Execution order: < P2 P0 P1 >

========================================
  All 3 processes finished successfully.
  Total steps taken: 3
  Final execution order: < P2 P0 P1 >
========================================
\end{lstlisting}
\end{outputbox}

\vfill
\begin{center}
{\color{rulegray}\rule{0.4\linewidth}{0.6pt}}\\[0.3cm]
{\small\textsc{End of Report}}
\end{center}

\end{document}
